% Part: incompleteness
% Chapter: arithmetization-syntax
% Section: coding-symbols

\documentclass[../../../include/open-logic-section]{subfiles}

\begin{document}

\olfileid{inc}{art}{cod}
\olsection{编码符号}

一阶逻辑的基本语言~$\Lang L$ 使用符号
\[
\lfalse \quad \lnot \quad \lor \quad \land \quad \lif \quad \lforall
\quad \lexists \quad \eq \quad ( \quad ) \quad ,
\]
以及可数的变元集与常项符号集，和任意元数的可数函数符号集与谓词符号集。
我们可以为每个符号分配一个唯一的数字作为其\emph{代码}，且不同符号的代码不同。
我们知道这是可行的，因为所有符号的集合是可数的，从而它与自然数集之间存在双射。
但我们还希望确保能够从代码中\emph{可计算地}还原出符号（以及关于它的一些信息，例如函数符号的元数）。
当然，实现这一点有很多方式。这里给出一种基于原始递归函数的方法。
（回忆一下，$\tuple{n_0, \dots, n_k}$ 是编码数序列 $n_0$, \dots, $n_k$ 的数。）

\begin{defn}
若 $s$ 是语言~$\Lang L$ 的符号，则其\emph{符号码}~$\scode s$ 定义如下：
\begin{enumerate}
\item 若 $s$ 是逻辑符号之一，则 $\scode s$ 由下表给出：
\[
\begin{array}{cccccccccc}
  \lfalse & \lnot & \lor & \land & \lif & \lforall \\
\tuple{0, 0} & \tuple{0, 1} & \tuple{0, 2} & \tuple{0, 3} &
\tuple{0, 4} & \tuple{0, 5} \\
\lexists & \eq & ( & ) & ,\\
\tuple{0, 6} & \tuple{0, 7} &
\tuple{0, 8} & \tuple{0, 9} & \tuple{0, 10}
\end{array}
\]
\item 若 $s$ 是第 $i$ 个变元 $\Obj v_i$，则 $\scode s = \tuple{1, i}$。
\item 若 $s$ 是第 $i$ 个常项符号~$\Obj c_i$，则
  $\scode s = \tuple{2, i}$。
\item 若 $s$ 是第 $i$ 个 $n$ 元函数符号~$\Obj f_i^n$，则
  $\scode s = \tuple{3, n, i}$。
\item 若 $s$ 是第 $i$ 个 $n$ 元谓词符号~$\Obj P_i^n$，则
  $\scode s = \tuple{4, n, i}$。
\end{enumerate}
\end{defn}

\begin{prop}
下列关系是原始递归的：
\begin{enumerate}
\item $\fn{Fn}(x, n)$ 成立当且仅当 $x$ 是某个~$i$ 对应的 $\Obj f^n_i$ 的代码，即 $x$ 是一个 $n$ 元函数符号的代码。
\item $\fn{Pred}(x, n)$ 成立当且仅当 $x$ 是某个~$i$ 对应的 $\Obj P^n_i$ 的代码，或者 $x$ 是 $\eq$ 的代码且 $n = 2$，即 $x$ 是一个 $n$ 元谓词符号的代码。
\end{enumerate}
\end{prop}

\begin{defn}
如果 $s_0, \dots, s_{n-1}$ 是一个符号序列，则其 \emph{Gödel数} 为 $\tuple{\scode{s_0}, \dots, \scode{s_{n-1}}}$。
\end{defn}

\begin{explain}
注意，\emph{代码}和\emph{Gödel数}是不同的。
例如，变元~$\Obj v_5$ 有一个代码~$\scode{\Obj
  v_5} = \tuple{1, 5} = 2^2\cdot 3^6$。
但是，变元~$\Obj v_5$ 作为一项，也是一个（长度为~$1$ 的）符号序列。
\emph{项}~$\Obj v_5$ 的\emph{Gödel数}~$\Gn{\Obj v_5}$ 是 $\tuple{\scode{\Obj v_5}} = 2^{\scode{\Obj
    v_5} + 1} = 2^{2^2\cdot 3^6 + 1}$。
\end{explain}

\begin{ex}
回想一下，如果 $k_0$, \dots, $k_{n-1}$ 是一个数序列，那么在素数幂编码下，该序列 $\tuple{k_0, \dots, k_{n-1}}$ 的代码是
\[
2^{k_0+1}\cdot3^{k_1+1}\cdot \dots \cdot p_{n-1}^{k_{n-1}+1},
\]
其中 $p_i$ 是第 $i$ 个素数（从 $p_0 = 2$ 开始）。
因此，举例来说，公式 $\eq[\Obj v_0][\Obj 0]$，或者更明确地写为 ${\eq}(\Obj v_0,\Obj c_0)$，其Gödel数是
\[
\tuple{\scode{\eq},\scode{(},\scode{\Obj v_0},\scode{,}, \scode{\Obj
    c_0},\scode{)}}.
\]
这里，$\scode{\eq}$ 是 $\tuple{0,7} = 2^{0+1}\cdot
3^{7+1}$，$\scode{\Obj v_0}$ 是 $\tuple{1,0} = 2^{1+1}\cdot3^{0+1}$，
依此类推。所以 $\Gn{=(\Obj v_0,\Obj c_0)}$ 是
\begin{multline*}
2^{\scode{=} + 1}\cdot 3^{\scode{(}+1}\cdot 5^{\scode{\Obj v_0}+1}
\cdot 7^{\scode{,} + 1} \cdot 11^{\scode{\Obj c_0}+1} \cdot
13^{\scode{)}+1} = \\
2^{2^1\cdot 3^8 + 1}\cdot 3^{2^1\cdot 3^9+1}\cdot 5^{2^2\cdot 3^1+1}
\cdot 7^{2^1\cdot 3^{11} + 1} \cdot 11^{2^3\cdot3^1+1} \cdot
13^{2^1\cdot3^{10}+1} = \\
2^{13\,123}\cdot 3^{39\,367}\cdot 5^{13}\cdot 7^{354\,295}\cdot11^{25}\cdot13^{118\,099}.
\end{multline*}
\end{ex}

\end{document}